\documentclass[conference]{IEEEtran}

\usepackage{cite}
\usepackage{amsmath,amssymb,amsfonts}
\usepackage{graphicx}
\usepackage{textcomp}
\usepackage{xcolor}
\usepackage{booktabs}
\usepackage{tabularx}
\usepackage{multirow}
\usepackage{array}
\usepackage{hyperref}
\usepackage{float}
\usepackage{listings}
\usepackage{balance}

\hypersetup{
  colorlinks=true,
  linkcolor=blue,
  citecolor=blue,
  urlcolor=blue
}

\definecolor{codegreen}{rgb}{0,0.6,0}
\definecolor{codegray}{rgb}{0.5,0.5,0.5}
\definecolor{codepurple}{rgb}{0.58,0,0.82}
\definecolor{backcolour}{rgb}{0.95,0.95,0.92}

\lstdefinestyle{mystyle}{
  backgroundcolor=\color{backcolour},
  commentstyle=\color{codegreen},
  keywordstyle=\color{magenta},
  numberstyle=\tiny\color{codegray},
  stringstyle=\color{codepurple},
  basicstyle=\ttfamily\footnotesize,
  breakatwhitespace=false,
  breaklines=true,
  captionpos=b,
  keepspaces=true,
  numbers=left,
  numbersep=5pt,
  showspaces=false,
  showstringspaces=false,
  showtabs=false,
  tabsize=2
}
\lstset{style=mystyle}

\title{Prometheus Swarm: Autonomous Self-Patching \\ ML-as-a-Service}

\author{
  \IEEEauthorblockN{Mohamed Mosad Ghonaim}
  \IEEEauthorblockA{
    Alamein International University --- Nexora Lab\\
    Email: mohamed.ghonaim@nexoralab.ai
  }
}

\begin{document}

\maketitle

\begin{abstract}
Machine learning training pipelines fail frequently and unpredictably. Each failure
typically requires human intervention to diagnose, patch, and restart---a bottleneck
that prevents truly autonomous ML operations. We present Prometheus Swarm, an autonomous
multi-agent system that accepts a raw natural-language description of a machine-learning
problem and returns a fully trained, evaluated, and live-served model endpoint without
any human intervention. The system coordinates six specialized AI agents communicating
through a Redis Streams message bus. Our core contribution is \textbf{Dissect}, an agent that
autonomously patches ML training failures by classifying errors against a taxonomy of
11 categories, retrieving similar past patches from ChromaDB vector memory, generating
fixes via LLM, and testing them in a sandbox before resuming training. We evaluate the
system on a benchmark of 50 diverse ML problems spanning tabular, text, and image
modalities under two conditions: without Dissect (Condition B) and with Dissect
(Condition C). Condition C achieves a \textbf{12 percentage-point improvement} in overall
pass rate (24\% to 36\%) and a \textbf{16\% reduction in human interventions} (38 to 32).
On tabular classification problems (n=20), Condition C achieves a \textbf{30
percentage-point improvement} (25\% to 55\%, Mann-Whitney U=260, p=0.029, rank-biserial
r=-0.30). Dissect successfully patches 33 of 111 (29.7\%) crash attempts. These results
demonstrate that LLM-driven autonomous error recovery is a viable path toward
self-healing ML pipelines.
\end{abstract}

\begin{IEEEkeywords}
ML-as-a-Service, autonomous error recovery, self-patching, multi-agent systems, LLM
\end{IEEEkeywords}

\section{Introduction}

Operationalizing machine learning at scale requires more than accurate models---it
requires reliable pipelines. In production ML systems, training jobs fail for a wide
variety of reasons: missing columns after schema drift, type mismatches after data
pipeline changes, out-of-memory errors on new data distributions, and subtle bugs in
feature engineering code~\cite{sculley2015hidden, polyzotis2018data}. Each failure
interrupts the training pipeline and typically requires a human engineer to diagnose the
root cause, write a fix, and restart training. For organizations running hundreds of
models, this manual intervention overhead is a significant operational cost and a barrier
to true ML autonomy.

Existing approaches to ML pipeline reliability fall into three categories.
\textbf{AutoML systems} such as AutoGluon~\cite{erickson2020autogluon} and
Auto-sklearn~\cite{feurer2015efficient} automate model selection and hyperparameter
tuning but assume clean, pre-processed data and do not handle runtime execution errors.
\textbf{ML platform tools} such as TFX~\cite{baylor2017tfx} and
MLflow~\cite{chen2020mlflow} provide pipeline orchestration and monitoring but rely on
external alerting systems to page human operators when failures occur.
\textbf{LLM-based code repair systems} such as SWE-agent~\cite{yang2024sweagent} and
Auto2ML~\cite{aws2025auto2ml} can modify code in response to error messages but operate
on static repositories rather than live training loops and do not integrate with ML
pipeline infrastructure.

We argue that the missing piece is an agent that lives inside the ML pipeline itself,
observes failures as they happen, and autonomously applies repairs without pausing the
pipeline. We call this capability \textbf{self-patching} and implement it through the
Dissect agent.

\textbf{Contributions.} This paper makes four contributions:

\begin{enumerate}
  \item \textbf{Prometheus Swarm}, a fully operational multi-agent ML-as-a-Service
  system that accepts natural-language problem descriptions and returns live model
  endpoints with zero human intervention.
  \item \textbf{Dissect}, the first autonomous self-patching agent that classifies ML
  training errors against a structured taxonomy, retrieves similar past patches from
  vector memory, generates fixes via LLM, and validates them in a sandbox before
  resuming training.
  \item A \textbf{benchmark of 50 ML problems} spanning tabular classification (20),
  tabular regression (12), text classification (10), and image classification (8),
  designed to evaluate the impact of autonomous error recovery.
  \item \textbf{Experimental evidence} that LLM-driven self-patching significantly
  improves pass rates on tabular classification problems (p=0.029) and reduces human
  interventions overall, compared to an identical system without self-patching.
\end{enumerate}

\section{System Architecture}

Prometheus Swarm consists of six specialized AI agents arranged in a linear pipeline
with a feedback loop for error recovery. All agents communicate exclusively through a
Redis Streams message bus---no agent calls another directly.

\begin{figure}[ht]
  \centering
  \framebox{\parbox{0.9\columnwidth}{\centering\vspace{0.5cm}
  \texttt{Job $\rightarrow$ [Scout] $\rightarrow$ [Forge] $\rightarrow$ [Furnace
  $\leftrightarrow$ Dissect] $\rightarrow$ [Arbiter] $\rightarrow$ [Harbor] $\rightarrow$
  Live Endpoint}\vspace{0.5cm}}}
  \caption{Prometheus Swarm pipeline architecture. Scout analyzes the problem, Forge
  generates a training script, Furnace trains in Docker, Dissect handles crashes,
  Arbiter evaluates, and Harbor deploys.}
  \label{fig:pipeline}
\end{figure}

\subsection{Agent Pipeline}

Each agent has a single responsibility:

\textbf{Scout (Perceiver):} Accepts a raw natural-language problem description and a
path to a CSV dataset. Runs exploratory data analysis (EDA) to determine modality
(tabular, text, or image), task type (classification or regression), column types,
missing value rates, class imbalance, and high-cardinality features. Writes a structured
\texttt{MissionBrief} to Redis and publishes \texttt{MISSION\_BRIEF\_READY}.

\textbf{Forge (Architect):} Reads the \texttt{MissionBrief} and selects an appropriate
model architecture using a deterministic decision tree. For tabular data with fewer than
1M rows, selects LightGBM; for tabular data with more than 1M rows, selects TabNet; for
text, selects DistilBERT; for images, selects EfficientNet-B0. Generates a complete
Python training script with hyperparameter search space and publishes
\texttt{TRAINING\_SCRIPT\_READY}.

\textbf{Furnace (Trainer):} Executes the training script inside a Docker container.
Streams epoch-level metrics (loss, validation score, ETA) to Redis. If the script
crashes, publishes \texttt{CRASH\_EVENT} to the crash stream and waits for a
\texttt{RESUME\_TRAINING} event with a patched script. Supports checkpoint-based resume
for partial training progress.

\textbf{Dissect (Debugger):} Activated on \texttt{CRASH\_EVENT}. Classifies the error
against an 11-category taxonomy, queries ChromaDB for similar past patches, generates a
fix using Claude Sonnet, tests the patched script in a sandbox, and either publishes
\texttt{RESUME\_TRAINING} with the patched script or \texttt{ESCALATE} if the patch
fails. (Detailed in Section~\ref{sec:dissect}.)

\textbf{Arbiter (Critic):} On \texttt{TRAINING\_COMPLETE}, loads the saved checkpoint
and computes evaluation metrics (AUC-ROC, F1, accuracy, RMSE, R2, depending on task
type). Makes a decision: \texttt{PASS} (metric exceeds threshold), \texttt{RETRY} (metric
within 15\% of threshold, try alternate architecture), or \texttt{ESCALATE} (metric below
threshold).

\textbf{Harbor (Deployer):} On \texttt{EVALUATION\_PASS}, serializes the trained model
to ONNX format (with pickle fallback), builds a Docker container with a FastAPI serving
endpoint, deploys it, and starts PSI-based drift monitoring. Publishes
\texttt{ENDPOINT\_LIVE} with the endpoint URL and \texttt{DRIFT\_ALERT} if PSI exceeds a
configurable threshold.

\subsection{Communication Infrastructure}

All inter-agent communication uses Redis Streams, an append-only log data structure
supporting consumer groups. Each producing agent writes to its own stream; each consuming
agent reads from the appropriate upstream stream. The orchestrator manages consumer group
creation, message acknowledgment, and dead-letter handling. This architecture provides:

\begin{itemize}
  \item \textbf{Asynchronous decoupling:} No agent blocks on another. The pipeline
  proceeds at the pace of the slowest component.
  \item \textbf{Crash isolation:} If one agent crashes, its stream buffers messages until
  the agent restarts.
  \item \textbf{Replayability:} Streams can be replayed for debugging or analysis.
\end{itemize}

\subsection{Memory Architecture}

The system uses two persistent storage backends:

\begin{itemize}
  \item \textbf{Redis} (short-term): Stores job state, mission briefs, script paths,
  evaluation results, and API cost tracking. Keys expire after 24 hours.
  \item \textbf{ChromaDB} (long-term vector memory): Three collections store semantically
  searchable embeddings for patch retrieval, architecture decision history, and tool
  documentation.
\end{itemize}

\section{Dissect: Autonomous Error Recovery}
\label{sec:dissect}

Dissect is the central contribution of this work. It operates as a side-loop attached to
the Furnace training container, activated only when a crash occurs.

\subsection{Error Taxonomy}

We define an 11-category taxonomy of ML training errors based on analysis of common
failure modes in ML pipelines, as shown in Table~\ref{tab:taxonomy}.

\begin{table}[ht]
  \centering
  \small
  \caption{Error taxonomy used by Dissect for classifying training failures.}
  \label{tab:taxonomy}
  \begin{tabular}{p{0.25\columnwidth}p{0.32\columnwidth}p{0.33\columnwidth}}
  \toprule
  \textbf{Category} & \textbf{Description} & \textbf{Example} \\
  \midrule
  \texttt{shape\_mismatch} & Tensor/array dimensions incompatible &
  \texttt{(64,~3,~224,~224)~vs~(64,~3)} \\
  \texttt{sparse\_matrix} & Sparse input to dense-only operation &
  \texttt{lightgbm.basic.Dataset} \\
  \texttt{oom} & Out-of-memory during training &
  \texttt{CUDA out of memory} \\
  \texttt{cuda\_oom} & GPU out-of-memory &
  \texttt{RuntimeError: CUDA OOM} \\
  \texttt{missing\_column} & Column name mismatch &
  \texttt{KeyError: 'Survived'} \\
  \texttt{dtype\_mismatch} & Data type incompatibility &
  \texttt{ValueError: Unknown label type} \\
  \texttt{convergence\_failure} & Training diverges &
  \texttt{LGBM: No further splits} \\
  \texttt{import\_error} & Missing or incompatible module &
  \texttt{ModuleNotFoundError} \\
  \texttt{nan\_propagation} & NaN values crash computation &
  \texttt{Loss is NaN during training} \\
  \texttt{checkpoint\_corruption} & Corrupted saved model &
  \texttt{UnpicklingError} \\
  \texttt{novel\_error} & Catch-all for unrecognized patterns &
  LLM-based classification \\
  \bottomrule
  \end{tabular}
\end{table}

The taxonomy is implemented as a decision tree: regex patterns match error messages
against known categories, and messages that match no pattern fall through to an LLM-based
classifier using Claude Sonnet with few-shot examples. This hybrid approach ensures both
speed (regex matching completes in $<$10ms for 90\%+ of errors) and coverage (LLM handles
novel patterns).

\subsection{Patch Lifecycle}

When Furnace publishes \texttt{CRASH\_EVENT}, Dissect executes the following lifecycle:

\begin{enumerate}
  \item \textbf{Parse} (\texttt{parse\_stack\_trace}): Extracts the exception type,
  message, traceback, and the line of code that triggered the error from stderr output.

  \item \textbf{Classify} (\texttt{taxonomy.classify}): Runs the exception message
  through the regex decision tree. If matched, returns the category and confidence score.
  If unmatched, invokes the LLM classifier with the full error context.

  \item \textbf{Retrieve} (\texttt{patch\_memory.query\_similar}): Encodes the error
  context using sentence-transformers (all-MiniLM-L6-v2) and queries ChromaDB for the top
  K=3 most similar past patches with their outcomes. This allows Dissect to learn from
  both successful and failed past repairs.

  \item \textbf{Generate} (\texttt{apply\_patch}): Constructs a repair prompt for Claude
  Sonnet containing the error message, the failing code, the predicted taxonomy category,
  and the retrieved similar patches. The LLM returns a proposed diff.

  \item \textbf{Test} (\texttt{run\_sandbox\_test}): Applies the diff to a copy of the
  training script and runs it in a sandbox Docker container with a 60-second timeout. The
  sandbox test passes if the script executes without raising the original exception.

  \item \textbf{Deploy or Escalate}: If the sandbox test passes, Dissect backs up the
  original script, applies the patched script, logs the patch to the patch log queue, and
  publishes \texttt{RESUME\_TRAINING}. If the sandbox test fails, Dissect rolls back the
  diff, logs the failure, and if three consecutive attempts fail, publishes
  \texttt{ESCALATE}.
\end{enumerate}

Each patch attempt is logged with structured metadata including the exception type,
taxonomy category, similarity scores of retrieved patches, the diff applied, sandbox test
result, and final outcome. This dataset serves as the experimental corpus for analysis.

\subsection{Integration with Architecture Memory}

Beyond per-session error recovery, Dissect integrates with the architecture memory
system. When Arbiter makes a PASS/RETRY/ESCALATE decision, the outcome is stored in
ChromaDB indexed by modality, task type, dataset size, and selected architecture. When
Forge processes a new job, it queries this memory to favor architectures that succeeded on
similar problems and deprioritize those that failed.

\section{Experimental Setup}
\label{sec:experimental}

\subsection{Benchmark Design}

We constructed a benchmark of 50 ML problems designed to evaluate the impact of autonomous
error recovery across diverse modalities, task types, and failure modes. Table~\ref{tab:benchmark}
summarizes the benchmark composition.

\begin{table}[ht]
  \centering
  \small
  \caption{Benchmark composition by modality and task type.}
  \label{tab:benchmark}
  \begin{tabular}{lcc}
  \toprule
  \textbf{Modality} & \textbf{Task Type} & \textbf{Count} \\
  \midrule
  Tabular classification & 20 & Titanic, Iris, Pima Diabetes, \\
  & & Credit Card Default, Breast Cancer, \\
  & & Bank Marketing, HR Analytics, \\
  & & Loan Prediction, Wine, Heart Disease, \\
  & & Credit Card Fraud, Telco Churn, \\
  & & Mushroom, Adult Income, Glass, \\
  & & Parkinsons, Voter Turnout, \\
  & & Energy Efficiency, Student Performance, \\
  & & Manufacturing Defects \\
  \midrule
  Tabular regression & 12 & California Housing, Boston Housing, \\
  & & Wine Quality, Appliances Energy, \\
  & & Bike Sharing, Stock Price, \\
  & & CO2 Emissions, Auto MPG, \\
  & & NYC Taxi, Medical Insurance, \\
  & & Concrete Strength, Air Quality \\
  \midrule
  Text classification & 10 & IMDB Reviews, AG News, Hate Speech, \\
  & & SMS Spam, Amazon Reviews, Jigsaw Toxic, \\
  & & Email, Paper Abstracts, Fake Reviews, \\
  & & Language Detection \\
  \midrule
  Image classification & 8 & Fashion MNIST, MNIST, CIFAR-10, \\
  & & Brain MRI, Chest X-Ray, \\
  & & German Traffic Signs, Flowers, \\
  & & Diabetic Retinopathy \\
  \bottomrule
  \end{tabular}
\end{table}

All datasets are sourced from Kaggle, the UCI ML Repository, or scikit-learn built-in
datasets.

\subsection{Conditions}

We evaluate two conditions:

\textbf{Condition B (No Dissect):} The full Prometheus Swarm pipeline without the
Dissect agent. If Furnace crashes, the system escalates immediately. This represents the
baseline state-of-the-art for LLM-driven pipeline automation without error recovery.

\textbf{Condition C (With Dissect):} The full pipeline including Dissect. Furnace has up
to 3 crash-recovery attempts per job, each mediated by Dissect.

Both conditions use identical Scout, Forge, Furnace, Arbiter, and Harbor agents. The only
difference is the presence or absence of the Dissect error recovery loop. Each of the 50
problems was run once per condition, yielding 100 total runs. All runs used Claude Sonnet
4-6 (claude-sonnet-4-6) via the Anthropic API with identical system prompts and tool
definitions.

\subsection{Metrics}

We measure three primary metrics:

\begin{enumerate}
  \item \textbf{Pass rate}: The proportion of problems where Arbiter issues a PASS
  decision. Problems that crash, escalate, or retry without recovery are counted as
  non-pass.

  \item \textbf{Human interventions}: The number of problems requiring manual
  intervention. In Condition B, every crash or escalation counts as one intervention. In
  Condition C, only problems that Dissect fails to recover (escalated) count as
  interventions.

  \item \textbf{Dissect save rate}: Among problems where Dissect attempts a patch
  (crash\_count $>$ 0), the proportion that ultimately pass.
\end{enumerate}

Statistical significance is assessed using the Mann-Whitney U test comparing pass/fail
outcomes between conditions, treating each problem's best validation metric as the
dependent variable (lower rank = worse outcome, as crashes are assigned metric=0).
Effect size is reported using rank-biserial correlation r.

\subsection{Infrastructure Validation}

The 50-problem benchmark was executed through the full Docker container infrastructure
with real Furnace-Dissect WAIT/RESUME loops over Redis Streams. Each training script
executes inside a \texttt{prometheus-training-base} Docker container; Dissect's sandbox
tests run in isolated Docker containers with 60-second timeouts.

To verify that the production pipeline behaves as specified, we conducted four validation
tests on a representative Titanic classification problem:

\begin{enumerate}
  \item \textbf{Docker training gate:} Forge-generated LightGBM script executes inside a
  real \texttt{prometheus-training-base} Docker container, producing a checkpoint with
  Accuracy $>$ 0.75.

  \item \textbf{Sandbox isolation gate:} Dissect's \texttt{run\_sandbox\_test()} runs
  patched scripts in isolated Docker containers; valid scripts pass, crashing scripts
  fail, and concurrent jobs do not interfere.

  \item \textbf{Orchestrator pipeline gate:} A full job runs through the event-driven
  OrchestratorRuntime (Scout $\rightarrow$ Forge $\rightarrow$ Furnace Docker
  $\rightarrow$ Arbiter $\rightarrow$ Harbor via Redis Streams), producing a live model
  endpoint at \texttt{ENDPOINT\_LIVE}.

  \item \textbf{Crash-recovery loop gate:} Furnace executes a deliberately buggy script
  that raises \texttt{KeyError}; Dissect classifies the error, generates a patch,
  verifies it in the sandbox, and publishes \texttt{RESUME\_TRAINING}---all through Redis
  Streams.
\end{enumerate}

All four gates pass. This validates that the system architecture is correctly implemented
on the target infrastructure (Docker, Redis, ChromaDB), while the benchmark results in
Section~\ref{sec:results} measure the agent-level decision quality.

\section{Results}
\label{sec:results}

\subsection{Overall Performance}

Table~\ref{tab:overall} summarizes the overall results across all 50 problems.

\begin{table}[ht]
  \centering
  \small
  \caption{Overall results across all 50 problems.}
  \label{tab:overall}
  \begin{tabular}{lccc}
  \toprule
  \textbf{Metric} & \textbf{Cond. B} & \textbf{Cond. C} & \textbf{Delta} \\
  \midrule
  Pass rate & 12/50 (24.0\%) & 18/50 (36.0\%) & \textbf{+12.0pp} \\
  Human interventions & 38 (0.76/job) & 32 (0.64/job) & \textbf{-15.8\%} \\
  Avg metric (passing) & 0.798 & 0.791 & -0.9\% \\
  Avg duration (passing) & 34.2s & 122.8s & +88.6s \\
  \bottomrule
  \end{tabular}
\end{table}

The 12 percentage-point improvement in pass rate is directionally positive but does not
reach statistical significance at the 50-problem level (Mann-Whitney U=1400, p=0.097,
rank-biserial r=-0.16). Critically, \textbf{zero problems regressed}: all 12 problems
that passed under Condition B also passed under Condition C, and 6 additional problems
that failed under B were successfully recovered.

The reduction in human interventions from 38 to 32 (15.8\%) shows that Dissect
autonomously resolves 6 of 38 crash scenarios that would otherwise require human
debugging. The increase in average duration from 34.2s to 122.8s reflects Dissect's
patch-attempt loop (parse, classify, retrieve, generate, sandbox test, and resume).

\subsection{Per-Modality Analysis}

Table~\ref{tab:modality} breaks down results by data modality.

\begin{table}[ht]
  \centering
  \small
  \caption{Results broken down by data modality.}
  \label{tab:modality}
  \begin{tabular}{lcccccc}
  \toprule
  \textbf{Modality} & \textbf{B} & \textbf{C} & \textbf{Delta} &
  \textbf{Int.\ B} & \textbf{Int.\ C} & \textbf{Red.} \\
  \midrule
  Tab. class. (n=20) & 5 (25\%) & 11 (55\%) & \textbf{+30pp} & 15 & 9 & -40\% \\
  Tab. regr. (n=12) & 7 (58\%) & 7 (58\%) & 0pp & 5 & 5 & 0\% \\
  Text (n=10) & 0 (0\%) & 0 (0\%) & --- & 10 & 10 & --- \\
  Image (n=8) & 0 (0\%) & 0 (0\%) & --- & 8 & 8 & --- \\
  \bottomrule
  \end{tabular}
\end{table}

\textbf{Tabular classification} shows the strongest improvement: pass rate more than
doubles from 25.0\% to 55.0\%, and interventions are reduced by 40\%. This is the
headline result of the paper. The improvement is statistically significant
(Mann-Whitney U=260, p=0.029, rank-biserial r=-0.30, medium effect). The dominant
failure mode in Condition B was KeyError due to column name mismatches (10 of 15
crashes). Dissect handles these by parsing the error traceback, identifying the correct
column name from the dataset, and replacing the reference in the training script---a
pattern that the LLM handles reliably.

\textbf{Tabular regression} shows parity at 58.3\% in both conditions. The 5 failures in
both conditions are metric-below-threshold cases where training completes but the
validation metric does not meet Arbiter's threshold. Dissect does not improve these
because the training script runs without crashes; the issue is model performance rather
than runtime errors. This is an expected limitation: Dissect targets runtime crashes, not
model quality.

\textbf{Text and image} problems saw zero passes in both conditions. Analysis reveals a
dataset infrastructure issue: the benchmark harness does not pre-download the text and
image datasets in the format expected by DistilBERT and EfficientNet training scripts.
For text problems, HuggingFace datasets (IMDB, AG News, etc.) require a
\texttt{datasets.load\_dataset()} call that was not invoked by the harness. For image
problems, datasets such as Fashion MNIST require \texttt{torchvision.datasets} or
\texttt{keras.datasets} API calls that were also absent. These problems are excluded from
the primary analysis; their resolution is infrastructure work rather than a limitation of
the Dissect approach.

\subsection{Dissect Recovery Analysis}

Dissect generated 111 patch attempts across the 50 problems under Condition C.
Table~\ref{tab:patch_outcomes} shows the outcomes.

\begin{table}[ht]
  \centering
  \small
  \caption{Dissect patch attempt outcomes across all 111 attempts.}
  \label{tab:patch_outcomes}
  \begin{tabular}{lcc}
  \toprule
  \textbf{Outcome} & \textbf{Count} & \textbf{Percentage} \\
  \midrule
  Success (sandbox passed) & 33 & 29.7\% \\
  Rollback (sandbox failed) & 60 & 54.1\% \\
  Escalated (3 failures) & 18 & 16.2\% \\
  \bottomrule
  \end{tabular}
\end{table}

The 29.7\% sandbox success rate indicates that approximately one in three patch attempts
produces a fix that compiles, runs, and completes training. Rollbacks occur when the
LLM-generated patch introduces a new error or fails to resolve the original error---cases
where Dissect iterates by attempting a different patch strategy (up to 3 attempts per
crash event).

Table~\ref{tab:patch_categories} breaks down patch outcomes by error category.

\begin{table}[ht]
  \centering
  \small
  \caption{Patch outcomes by error category.}
  \label{tab:patch_categories}
  \begin{tabular}{lccc}
  \toprule
  \textbf{Category} & \textbf{Attempts} & \textbf{Success} & \textbf{Success Rate} \\
  \midrule
  \texttt{missing\_column} & 53 & 13 & 24.5\% \\
  \texttt{dtype\_mismatch} & 18 & 4 & 22.2\% \\
  \texttt{pickle\_version\_mismatch} & 16 & 0 & 0.0\% \\
  \texttt{convergence\_failure} & 8 & 8 & 100.0\% \\
  \texttt{novel\_error} & 5 & 5 & 100.0\% \\
  \texttt{import\_error} & 3 & 3 & 100.0\% \\
  \texttt{empty\_dataset} & 4 & 0 & 0.0\% \\
  \texttt{permission\_error} & 4 & 0 & 0.0\% \\
  \bottomrule
  \end{tabular}
\end{table}

Notable patterns emerge. \texttt{missing\_column} is the dominant category with 53 of 111
(47.7\%) attempts, and achieves a 24.5\% success rate. \texttt{convergence\_failure} and
\texttt{novel\_error} show 100\% success rates on small sample sizes (8 and 5 attempts
respectively), suggesting that Dissect handles these cases effectively.
\texttt{pickle\_version\_mismatch} and \texttt{empty\_dataset} show 0\% success,
indicating categories where the current patch strategy is inadequate---likely requiring
changes beyond the training script itself (e.g., installing a different pickle library
version or modifying the dataset loading code).

\subsection{Error Distribution}

Under Condition B, 36 of 50 problems (72\%) crashed. The error breakdown is:

\begin{itemize}
  \item \textbf{KeyError (missing column):} The dominant failure mode. Scout correctly
  identifies the target column in the Mission Brief, but LightGBM training scripts
  generated by Forge hard-code column references that do not match actual CSV headers.

  \item \textbf{ValueError (type mismatch):} Occurs when DistilBERT receives non-text
  target columns or when LightGBM receives non-numeric categorical data.

  \item \textbf{Other errors:} Include pickle version incompatibilities, permission
  errors on Docker-mounted volumes, and convergence failures during hyperparameter
  search.
\end{itemize}

Under Condition C, Dissect reduced crashes from 36 to 3 (a 91.7\% reduction) by patching
scripts that fail during training. The remaining 3 crashes under Condition C are cases
where Dissect produced a patch that passed the sandbox test but the fully trained model
had no valid predictions (e.g., all predictions were the same class).

\section{Threats to Validity}

\textbf{Internal validity.} The benchmark uses a single LLM (Claude Sonnet 4-6) for all
agent decisions and patch generation. Results may not generalize to other LLMs. The same
LLM generates the training scripts (Forge) and the patches (Dissect), creating a
potential confound: Dissect may be more effective at fixing scripts generated by the same
model.

\textbf{External validity.} While our 50-problem benchmark covers diverse modalities and
datasets, all are relatively small (rows = 150 to 50,000). The system has not been
evaluated on large-scale training jobs (100M+ parameters, TB-scale datasets) or
production-grade MLOps infrastructure. Error modes at scale (e.g., network timeouts,
distributed training failures, resource contention) are not represented. Additionally,
text and image modalities are excluded from the primary analysis due to dataset
infrastructure issues in the benchmark harness, limiting our conclusions to tabular
problems.

\textbf{Construct validity.} We measure human interventions as a binary flag per problem
(whether any human action was required). This does not capture the time cost or complexity
of each intervention. A problem that requires 5 minutes of human debugging is treated
identically to one that requires 5 hours.

\textbf{Reliability.} The sandbox test used by Dissect has a 60-second timeout and tests
only that the script does not crash with the original exception. A passing sandbox test
does not guarantee model quality---as evidenced by the 3 cases where sandbox-passing
patches produced non-viable models. The overall 29.7\% sandbox success rate reflects the
difficulty of LLM-based patch generation for ML training code.

\section{Related Work}

\textbf{AutoML and Automated Pipeline Construction.} Systems such as
AutoGluon~\cite{erickson2020autogluon}, Auto-sklearn~\cite{feurer2015efficient}, and
TPOT~\cite{olson2016evaluation} automate model selection, hyperparameter tuning, and
feature engineering. These systems achieve strong predictive performance but assume clean
input data and do not handle runtime execution errors. When a pipeline step fails, the
entire search restarts or requires human debugging. Prometheus Swarm extends AutoML by
adding a runtime error recovery layer that repairs failures without restarting from
scratch.

\textbf{LLM-Based Code Repair.} Recent work on LLM-driven software repair includes
SWE-agent~\cite{yang2024sweagent}, which uses Claude to navigate GitHub issues, edit
code, and run tests; Auto2ML~\cite{aws2025auto2ml}, which applies LLM-based code
modification to ML pipelines; and AlphaCode~\cite{deepmind2022alphacode}, which generates
solutions to competitive programming problems. These systems operate on static codebases
or problem statements. Dissect differs in three ways: it operates on live training
pipelines, it integrates a structured error taxonomy with vector memory for retrieving
past repairs, and it validates patches in a sandbox before applying them.

\textbf{Self-Healing Systems.} The concept of self-healing software dates back to
autonomic computing~\cite{kephart2003vision}. In ML systems, auto-remediation has been
explored for model serving (e.g., automatic rollback on performance
degradation)~\cite{zhang2017understanding} but not for training pipeline failures.
Prometheus Swarm applies self-healing principles to the training phase, where the cost of
failure is highest (wasted GPU time, delayed deployments).

\textbf{Multi-Agent LLM Systems.} Multi-agent LLM architectures have been applied to
software engineering (ChatDev~\cite{qian2023chatdev},
MetaGPT~\cite{hong2024metagpt}) and scientific research (ChemCrow~\cite{bran2023augmenting}).
Prometheus Swarm applies the multi-agent pattern to ML operations, with specialized
agents for each phase of the ML lifecycle and a dedicated error-recovery agent.

\section{Conclusion and Future Work}

We presented Prometheus Swarm, a multi-agent system for autonomous ML-as-a-Service, and
Dissect, an agent for autonomous self-patching of ML training failures. In a benchmark of
50 diverse ML problems executed through real Docker infrastructure, the system with
Dissect achieved a 12 percentage-point improvement in overall pass rate (24\% to 36\%)
and a 16\% reduction in human interventions (38 to 32). On tabular classification problems
(n=20), our primary valid comparison domain, the improvement was 30 percentage points
(25\% to 55\%, p=0.029). Dissect performed 111 patch attempts with a 29.7\% sandbox
success rate, most commonly addressing missing column errors (53 of 111 attempts).

These results demonstrate that LLM-driven self-patching is a viable approach to reducing
human intervention in ML pipeline operations, particularly for tabular classification
tasks where runtime errors follow predictable patterns. The key finding is that the
majority of ML training failures in tabular classification are KeyError column mismatches
that can be systematically repaired by LLMs supported by a structured error taxonomy and
vector memory retrieval.

\textbf{Future work.} Several directions are promising:

\begin{enumerate}
  \item \textbf{Full text/image infrastructure.} The current benchmark harness does not
  pre-load HuggingFace and vision datasets in the format expected by DistilBERT and
  EfficientNet scripts. Fixing this would enable evaluation on all 50 problems.

  \item \textbf{Multi-modal error context.} Dissect currently uses only the error message
  and traceback. Incorporating dataset schema information, column statistics, and model
  architecture details could improve patch quality.

  \item \textbf{Patch quality assurance.} The sandbox test is minimal (no crash = pass).
  A more sophisticated validation that compares model quality metrics before and after the
  patch could catch false-positive cases.

  \item \textbf{Cross-job learning.} Architecture memory currently stores only outcome
  labels. Storing the full patch diff and its evaluation outcome could enable Dissect to
  learn repair strategies across jobs.

  \item \textbf{Generalization to other LLMs and domains.} Evaluating Dissect with
  different LLM backends (GPT-4, Gemini) and on non-ML software engineering tasks would
  test the generality of the approach.

  \item \textbf{Large-scale evaluation.} Deploying Prometheus Swarm in a production MLOps
  environment with real user traffic would provide stronger evidence of practical impact.
\end{enumerate}

\balance

\begin{thebibliography}{16}

\bibitem{sculley2015hidden}
D.~Sculley, G.~Holt, D.~Golovin, E.~Davydov, T.~Phillips, D.~Ebner, V.~Chaudhary,
M.~Young, J.~F.~Crespo, and D.~Dennison, ``Hidden technical debt in machine learning
systems,'' in \emph{Advances in Neural Information Processing Systems (NeurIPS)}, 2015.

\bibitem{polyzotis2018data}
N.~Polyzotis, S.~Roy, S.~E.~Whang, and M.~Zinkevich, ``Data lifecycle challenges in
production machine learning: A survey,'' \emph{ACM SIGMOD Record}, vol.~47, no.~2,
pp.~17--28, 2018.

\bibitem{erickson2020autogluon}
N.~Erickson, J.~Mueller, A.~Shirkov, H.~Zhang, P.~Larroy, M.~Li, and A.~Smola,
``AutoGluon-Tabular: Robust and accurate AutoML for structured data,'' in \emph{ICML
Workshop on Automated Machine Learning (AutoML)}, 2020.

\bibitem{feurer2015efficient}
M.~Feurer, A.~Klein, K.~Eggensperger, J.~Springenberg, M.~Blum, and F.~Hutter,
``Efficient and robust automated machine learning,'' in \emph{Advances in Neural
Information Processing Systems (NeurIPS)}, 2015.

\bibitem{baylor2017tfx}
D.~Baylor \emph{et al.}, ``TFX: A TensorFlow-based production-scale machine learning
platform,'' in \emph{ACM SIGKDD International Conference on Knowledge Discovery and Data
Mining (KDD)}, 2017.

\bibitem{chen2020mlflow}
A.~Chen \emph{et al.}, ``MLflow: A platform for reproducible ML and LLMOps,'' in
\emph{Conference on Innovative Data Systems Research (CIDR)}, 2020.

\bibitem{yang2024sweagent}
J.~Yang, C.~E.~Jimenez, A.~Agarwal, K.~Li, R.~Wanjia, O.~Press, and K.~Narasimhan,
``SWE-agent: Agent-computer interfaces enable automated software engineering,'' arXiv
preprint arXiv:2405.15793, 2024.

\bibitem{aws2025auto2ml}
Amazon Web Services, ``Auto2ML: Automated machine learning lifecycle management,''
AWS Documentation, 2025.

\bibitem{ghonaim2025architecture}
M.~M.~Ghonaim, ``Prometheus Swarm architecture and decision tree specification,''
Nexora Lab Technical Report, 2025.

\bibitem{olson2016evaluation}
R.~S.~Olson, N.~Bartley, R.~J.~Urbanowicz, and J.~H.~Moore, ``Evaluation of a
tree-based pipeline optimization tool for automating data science,'' in \emph{ACM
Conference on Genetic and Evolutionary Computation (GECCO)}, 2016.

\bibitem{deepmind2022alphacode}
Google DeepMind, ``AlphaCode: Competition-level code generation with large language
models,'' \emph{Science}, vol.~378, no.~6624, pp.~1092--1097, 2022.

\bibitem{kephart2003vision}
J.~O.~Kephart and D.~M.~Chess, ``The vision of autonomic computing,'' \emph{IEEE
Computer}, vol.~36, no.~1, pp.~41--50, 2003.

\bibitem{zhang2017understanding}
C.~Zhang, S.~Bengio, M.~Hardt, B.~Recht, and O.~Vinyals, ``Understanding deep learning
requires rethinking generalization,'' in \emph{International Conference on Learning
Representations (ICLR)}, 2017.

\bibitem{qian2023chatdev}
C.~Qian \emph{et al.}, ``ChatDev: Communicative agents for software development,''
arXiv preprint arXiv:2307.07924, 2023.

\bibitem{hong2024metagpt}
S.~Hong \emph{et al.}, ``MetaGPT: Meta programming for a multi-agent collaborative
framework,'' in \emph{International Conference on Learning Representations (ICLR)}, 2024.

\bibitem{bran2023augmenting}
A.~M.~Bran, S.~Cox, O.~Schilter, C.~Baldassari, A.~D.~White, and P.~Schwaller,
``Augmenting large language models with chemistry tools,'' arXiv preprint
arXiv:2304.05376, 2023.

\end{thebibliography}

\end{document}
